\documentclass[10pt]{article}
\usepackage[margin=12mm]{geometry}
\usepackage{graphicx}
\usepackage{caption}
\pagestyle{empty}
\captionsetup{font=small,labelfont=bf}

\begin{document}
\begin{figure}[p]
\centering
\includegraphics[width=\textwidth]{fig4_ce_subtree_map_panel.pdf}

\caption{\textbf{The \emph{C.\@ elegans} lineage contains local subtrees
whose natural terminal assignment is exactly Pareto-optimal.}
%
Each circle is a lineage subtree with at least 12 usable terminal cells,
connected as the natural lineage; the circle area is proportional to its
usable terminal-cell count, and the pie shows the fate
composition of those biological terminal cells (terminal cells without valid
tracking or expression data are excluded). A solid circumference marks a
subtree whose natural terminal assignment is exactly Pareto-optimal for at
least one objective weighting; a dashed circumference marks an off-front
assignment. Subtree-name and cell-count color indicates the maximum fraction
of natural edges retained by any assignment on that subtree's Pareto front,
using the same blue retention scale as Figures~2, 3, and 6.
%
All 16 on-front subtrees contained at most 35 usable terminal cells; exact
recovery was absent among larger investigated subtrees.}
\label{fig:ce_subtree_map}
\end{figure}
\end{document}
